%Prva naredba uvijek definira klasu dokumenta. Mi koristimo 'article'

\documentclass[a4paper,12pt,oneside]{article}



%***********************************************************************************
% Sada slijedi niz poziva Paketa funkcija cije ce se naredbe kasnije koristiti:
%***********************************************************************************

%Paket za upravljanje sa listama
\usepackage{enumitem}

% Paket za podesavanje brojeva stranica
\usepackage{scrlayer-scrpage}                  
	\ifoot[]{}
	\cfoot[]{}
	\ofoot[\pagemark]{\pagemark}
	\pagestyle{scrplain}
	
	
% Najcesce koristeni paketi za matematiku i formatiranje rada:
\usepackage{xurl}
\usepackage{url}
\usepackage{amsmath, amsthm}
\usepackage{mathtools}
\usepackage{latexsym,amsfonts,amssymb,amsbsy,graphics}
\usepackage{graphicx}
\usepackage[croatian]{babel} %Omogucava pisanje hrvatskih znakoba
\usepackage{gensymb}
\usepackage{pdfpages}  %Omogucava importiranje stranica pdf fila u dokument
\usepackage{pdflscape} %Omogucava okretanje pojedine stranice landscape dok su ostale portrait


% Paket za postavke naslova i podnaslova
\usepackage{titlesec}  
	\newcommand{\periodafter}[1]{#1.}
	\titleformat{\section}[hang]{\bf\large\uppercase}{\thesection.}{5pt}{}{}
	\titleformat{\subsection}[hang]{\bf\large}{\thesubsection.}{3pt}{}{}

% Dva paketa za imenovanje slika i tablica
\usepackage{subcaption}
\usepackage[labelfont=bf,textfont=it,labelsep = colon]{caption}  
	\captionsetup[figure]{name=SL.}
	\captionsetup[table]{name=Tab.}

% Paket sa bojama za formatiranje svega gdje se mjenja boja
\usepackage{xcolor}

% Paket za omogucavanje importiranja programskih kodova u tekst
\usepackage{listings} 
	\lstset{
  		numbers=left,
  		numberstyle=\tiny,
  		columns=fullflexible,
  		breaklines=true,
  		postbreak=\mbox{\textcolor{red}{$\hookrightarrow$}\space},
  		xleftmargin=5.0ex
	}
	\renewcommand{\lstlistingname}{Primjer koda}% Listing -> Algorithm

% Paket za formatiranje Sadrzaja
\usepackage[dotinlabels]{titletoc}
	\dottedcontents{section}[5em]{\bfseries}{2em}{1pc}
	\dottedcontents{subsection}[7em]{}{2em}{1pc}
	
	
	
% Najbolji paket u svemiru za matematicko pisanje tekstova
\usepackage{physics}

% Paket za postavljanje linkova na koje se moze klikat u pdfu (prema formuli broj x- click odvede na to)
\usepackage[colorlinks,linkcolor={red!50!black},citecolor={blue!50!black},urlcolor={blue!80!black}]{hyperref}


% Postavke margina stranica - NE DIRAJ

\setlength{\topmargin}{-0.7in}             
\setlength{\oddsidemargin}{-2mm}
\setlength{\evensidemargin}{-2mm}
\setlength{\textwidth}{170mm}
\setlength{\textheight}{256mm}
\setlength{\parindent}{1cm}


%**************************************************************************************************************
% Prostor za tvoje pakete je ovdje:
%**************************************************************************************************************

\usepackage{booktabs}
\usepackage{multirow}
\usepackage{array}
\usepackage{float}


%**************************************************************************************************************




%Kraj preamble naredbi, krece se sa pisanjem dokumenta
%**************************************************************************************************************


\begin{document}

%**************************************************************************************************************
% Postavke za numeriranje, prorede i ostalo
%**************************************************************************************************************

\setcounter{MaxMatrixCols}{20}                          %Postavka da matrica moze imati vise od 10 stupaca
\renewcommand{\baselinestretch}{1.5}                    %postavke za prored 1.5%
\renewcommand{\thefigure}{\thesection.\arabic{figure}}  %postavke za imenovanje slika po poglavljima%
\renewcommand{\thetable}{\thesection.\arabic{table}}   %postavke za imenovanje tablica po poglavljima%
\renewcommand{\theequation}{\thesection-\arabic{equation}} %postavke za brojanje jednadzbi po poglavljima%
\renewcommand{\refname}{LITERATURA}						% Uppercase Literatura

%**************************************************************************************************************




%**************************************************************************************************************
% NASLOVNA STRANICA
%**************************************************************************************************************

% ============================================================
% 00-naslovnica.tex --- Prema FERIT predlo\v{s}ku
% ============================================================

\begin{titlepage}

\begin{center}
{\sc\Large sveu\v{c}ili\v{s}te josipa jurja strossmayera u osijeku}\\
{\sc\Large fakultet elektrotehnike, ra\v{c}unarstva i informacijskih tehnologija osijek}\\
\bigskip

\vspace*{2cm}

{\large Sveu\v{c}ili\v{s}ni diplomski studij Ra\v{c}unarstvo}\\

\vspace*{5cm}

{\sc\LARGE analiza algoritama zasnovanih na smote za\\ rukovanje problemom neuravnote\v{z}enosti klasa}

\vspace*{1cm}

{\large Diplomski rad}\\

\vspace*{3cm}

{\Large Kristian D\v{z}oi\'{c}}\\

\vspace*{8cm}

{\normalsize Osijek, 2026.}
\end{center}

\end{titlepage}




%**************************************************************************************************************
% OBRASCI I FORMULARI
%**************************************************************************************************************

% Ubaciti pdf verzije D1 obrazca i izjave o originalnosti. 
% Sustav MAK automatski ume\'ce ove dokumente.


%\newpage
%\includepdf{D1_obrazac.pdf}
%\newpage
%\includepdf{Originalnost.pdf}
%\newpage




%**************************************************************************************************************
% TABLE OF CONTENTS
%**************************************************************************************************************

{
	\hypersetup{linkcolor=black}
	\tableofcontents
}
%**************************************************************************************************************



\newpage

% ================================================================
% 1. UVOD
% ================================================================
\section{UVOD}
% ============================================================
% 01-uvod.tex — PRVA VERZIJA
% ============================================================

Ukratko navesti \v{s}to je to klasifikacija i \v{s}to joj je cilj te
neke od primjena koje obuhva\'{c}a. Sa\v{z}eto navesti problem
neuravnote\v{z}enosti klasa (engl.\ class imbalance) i kako utje\v{c}e
na u\v{c}inkovitost klasifikatora, osvrnuti se na to kako je manjinska
klasa (engl.\ minority class) \v{c}esto upravo ona od interesa. Treba
spomenuti preuzorkovanje (engl.\ oversampling) i poduzorkovanje (engl.\
undersampling) kao na\v{c}ine rukovanja navedenim problemom koji se
koriste samo dostupnim podacima.

Nadalje, navesti algoritam SMOTE (Synthetic Minority Over-sampling
Technique) kao popularni i u\v{c}inkoviti pristup preuzorkovanju koji
generira sinteti\v{c}ke primjere manjinske klase. Zbog nedostataka
temeljnog SMOTE-a --- prekomjerne generalizacije, osjetljivosti na
\v{s}um, generiranja u podru\v{c}jima ve\'{c}inske klase --- predlo\v{z}ena
su brojna pro\v{s}irenja: Borderline-SMOTE, ADASYN, Safe-Level-SMOTE,
K-Means SMOTE, SVM-SMOTE, SMOTE-ENN, SMOTE-Tomek te geometrijske
varijante poput G-SMOTE i Random SMOTE. Svaka izvedenica nastoji
rije\v{s}iti odre\dj{}eni nedostatak originala --- bilo druga\v{c}ijim
na\v{c}inom odabira primjera za preuzorkovanje, bilo druga\v{c}ijim
na\v{c}inom stvaranja sinteti\v{c}kih uzoraka. Radi cjelovitosti,
u obzir su uzeti i alternativni pristupi: cost-sensitive u\v{c}enje,
ansambl metode te pristupi zasnovani na neuronskim mre\v{z}ama.

Treba naglasiti da ne postoji univerzalno najbolji algoritam za
preuzorkovanje --- u\v{c}inkovitost pojedine varijante ovisi o
karakteristikama skupa podataka: omjeru neuravnote\v{z}enosti,
dimenzionalnosti, prisutnosti \v{s}uma i geometrijskoj strukturi. Cilj
ovog rada jest eksperimentalno usporediti navedene algoritme na
skupovima podataka razli\v{c}itih karakteristika te statisti\v{c}ki
utvrditi zna\v{c}ajnost uo\v{c}enih razlika.

Ostatak rada organiziran je na sljede\'{c}i na\v{c}in. Drugo poglavlje
obra\dj{}uje problem neuravnote\v{z}enosti klasa i temeljni SMOTE
algoritam, uklju\v{c}uju\'{c}i mjere za evaluaciju i ote\v{z}avaju\'{c}e
faktore. Tre\'{c}e poglavlje donosi detaljan pregled izvedenica i
alternativnih pristupa. \v{C}etvrto poglavlje opisuje ostvareno
programsko rje\v{s}enje --- arhitekturu, implementirane algoritme,
na\v{c}in kori\v{s}tenja te web su\v{c}elje za vizualizaciju
generiranja sinteti\v{c}kih primjera. Peto poglavlje sadr\v{z}i
eksperimentalnu analizu, od opisa podataka i postavki eksperimenata
do statisti\v{c}ke obrade i diskusije rezultata.


% ================================================================
% 2. PROBLEM NEURAVNOTE\v{Z}ENOSTI KLASA I ALGORITAM SMOTE
% ================================================================
\section{PROBLEM NEURAVNOTE\v{Z}ENOSTI KLASA I ALGORITAM SMOTE}
% ============================================================
% 02-neuravnotezenost-i-smote.tex — PRVA VERZIJA
% ============================================================

\subsection{Klasifikacija i problem neuravnote\v{z}enosti klasa}

U ovom potpoglavlju bit \'{c}e opisano \v{s}to predstavlja klasifikacija
\cite{duda2000,theodoridis2008} i \v{s}to ima za cilj, ograni\v{c}avaju\'{c}i
se pritom na skupove podataka $\mathcal{A} \subset \mathbb{R}^m$ gdje su
sve zna\v{c}ajke koje opisuju podatke realne. Spomenut \'{c}e se neke
uobi\v{c}ajene primjene klasifikacije u razli\v{c}itim podru\v{c}jima,
uklju\v{c}uju\'{c}i medicinsku dijagnostiku \cite{odajima2014,kamel2019} i
druge domene.

Zatim se opisuje \v{s}to je neuravnote\v{z}enost klasa i na koji
na\v{c}in utje\v{c}e na problem klasifikacije --- u pravilu, rezultira
pristrano\v{s}\'{c}u klasifikatora prema ve\'{c}inskoj klasi
\cite{he2009,lopez2013}. Navest \'{c}e se za\v{s}to se radi o zna\v{c}ajnom
problemu i za\v{s}to je nu\v{z}no rukovati njime na odgovaraju\'{c}i
na\v{c}in, posebice jer je manjinska klasa \v{c}esto upravo ona od
interesa u stvarnim primjenama. Formalno \'{c}e se definirati omjer
neuravnote\v{z}enosti $IR = N_{\text{maj}} / N_{\text{min}}$ te
kategorizacija skupova prema vrijednosti IR-a.

Ukratko se opisuju poduzorkovanje i preuzorkovanje kao na\v{c}ini
rukovanja navedenim problemom. Iako nema univerzalno boljeg pristupa i
u\v{c}inkovitost ovisi o skupu podataka i klasifikatoru, preuzorkovanje
se \v{c}e\v{s}\'{c}e koristi jer kod poduzorkovanja postoji opasnost od
gubljenja bitnih uzoraka ve\'{c}inske klase \cite{bajer2019}. U ovom
dijelu ukomponirat \'{c}e se i podjela skupa podataka na podskup za
treniranje i testiranje.

U odnosu na postoje\'{c}u literaturu, dodatno \'{c}e se obratiti
pozornost na ote\v{z}avaju\'{c}e faktore koji prate problem
neuravnote\v{z}enosti: preklapanje klasa, unutar-klasni disbalans,
prisutnost \v{s}uma u podacima, mali apsolutni broj manjinskih primjera
te manifold struktura podataka \cite{krawczyk2016}.

\subsubsection{Neki popularni algoritmi za klasifikaciju}

Postoje brojni algoritmi za klasifikaciju
\cite{duda2000,theodoridis2008} --- nisu svi jednako prikladni za sve
probleme i odabir klasifikatora (engl.\ model selection) predstavlja
zaseban problem u strojnom u\v{c}enju. U nastavku se navodi osam
klasifikatora koji \'{c}e se koristiti u eksperimentalnom dijelu
rada, uz princip rada, klju\v{c}ne parametre te prednosti i nedostatke:

\begin{itemize}
    \item Stablo odluke (engl.\ Decision Tree, DT) --- rekurzivno dijeli
    podatke temeljem vrijednosti zna\v{c}ajki, sklono prekomjernom
    prilago\dj{}avanju.
    \item Slu\v{c}ajna \v{s}uma (engl.\ Random Forest, RF) --- ansambl
    stabala odluke s baggingom, smanjuje varijancu i robusnija je od
    pojedina\v{c}nog stabla.
    \item XGBoost (engl.\ Extreme Gradient Boosting) --- gradient boosting
    ansambl s regularizacijom, \v{c}esto posti\v{z}e vrhunske rezultate
    na tabli\v{c}nim podacima; ima ugra\dj{}enu podr\v{s}ku za
    neuravnote\v{z}ene klase putem parametra \texttt{scale\_pos\_weight}.
    \item Logisti\v{c}ka regresija (LR) --- linearni model s vjerojatnosnim
    izlazom, jednostavan i \v{c}esto dobar baseline.
    \item Stroj s potpornim vektorima (engl.\ Support Vector Machine, SVM)
    --- maksimizira marginu izme\dj{}u klasa, RBF jezgra omogu\'{c}uje
    nelinearnu separaciju.
    \item Algoritam $k$-najbli\v{z}ih susjeda ($k$-NN) --- klasificira
    temeljem klasa $k$ najbli\v{z}ih primjera, osjetljiv na distribuciju
    i skalu zna\v{c}ajki.
    \item Gaussov naivni Bayes (GNB) --- pretpostavlja normalnu distribuciju
    zna\v{c}ajki i me\dj{}usobnu nezavisnost, brz i jednostavan.
    \item Vi\v{s}eslojni perceptron (engl.\ Multi-Layer Perceptron, MLP)
    --- jednostavna neuronska mre\v{z}a s jednim skrivenim slojem,
    predstavlja osnovnu deep learning referentnu to\v{c}ku u usporedbi.
\end{itemize}

\subsubsection{Vrednovanje u\v{c}inkovitosti algoritama klasifikacije}

Objasnit \'{c}e se \v{s}to predstavlja matrica zbunjenosti (engl.\
confusion matrix) i kako se dobiva. Uobi\v{c}ajeno se u\v{c}inkovitost
klasifikatora vrednuje to\v{c}no\v{s}\'{c}u klasifikacije (engl.\
classification accuracy), no ona nije prikladna kod neuravnote\v{z}enih
klasa --- klasifikator mo\v{z}e posti\'{c}i visoku to\v{c}nost
predvi\dj{}aju\'{c}i isklju\v{c}ivo ve\'{c}insku klasu.

Zatim \'{c}e se opisati prikladnije mjere za neuravnote\v{z}ene podatke,
polaze\'{c}i od elemenata matrice zbunjenosti (TP, TN, FP, FN):
F-1 mjera kao harmonijska sredina preciznosti i odziva, G-Mean
kao geometrijska sredina odziva i specifi\v{c}nosti, AUC-ROC kao povr\v{s}ina
ispod ROC krivulje neovisna o pragu odluke, Balanced Accuracy kao
aritmeti\v{c}ka sredina odziva i specifi\v{c}nosti te Matthewsov
koeficijent korelacije (MCC) koji uzima u obzir sva \v{c}etiri elementa
matrice zbunjenosti. Za svaku mjeru dat \'{c}e se formula, interpretacija
i rasprava o prikladnosti. Ograni\v{c}it \'{c}e se na slu\v{c}aj binarne
klasifikacije gdje je pozitivna klasa uobi\v{c}ajeno manjinska.
Pritom \'{c}e se koristiti literatura \cite{sokolova2009,japkowicz2011,lopez2013,he2009}.

\subsection{Pregled pristupa rje\v{s}avanju problema}

U ovom potpoglavlju dat \'{c}e se krovni pregled \v{c}etiriju kategorija
pristupa rje\v{s}avanju problema neuravnote\v{z}enosti klasa, pri \v{c}emu
\'{c}e se detaljni opisi pojedinih metoda obraditi u tre\'{c}em poglavlju
\cite{fernandez2018book}.

Prvo, metode na razini podataka koje obuhva\'{c}aju
preuzorkovanje manjinske klase, poduzorkovanje ve\'{c}inske klase i
njihove hibridne kombinacije. Prednost ovih metoda je neovisnost o
klasifikatoru i jednostavnost primjene \cite{bajer2019,dudjak2020}.
Drugo, metode na razini algoritma koje uklju\v{c}uju cost-sensitive
u\v{c}enje i modifikacije klasifikatora \cite{elkan2001}. Tre\'{c}e,
ansambl metode koje kombiniraju resampliranje s bagging i boosting
pristupima \cite{galar2012}. \v{C}etvrto, deep learning pristupi koji
obuhva\'{c}aju modifikacije loss funkcija \cite{lin2017} i generativne
modele \cite{goodfellow2014}.

\subsection{Algoritam SMOTE}

U ovom potpoglavlju detaljno \'{c}e se opisati temeljni SMOTE algoritam
\cite{chawla2002}. Prvo \'{c}e se ukratko opisati nasumi\v{c}no
preuzorkovanje kao motivacija za nastanak SMOTE-a --- kod nasumi\v{c}nog
preuzorkovanja dupliciraju se postoje\'{c}i primjeri manjinske klase,
\v{s}to \v{c}esto dovodi do prekomjernog prilago\dj{}avanja. SMOTE ovaj
problem rje\v{s}ava generiranjem novih, sinteti\v{c}kih primjera.

Zatim \'{c}e se dati detaljan opis koraka algoritma: za svaki primjer
manjinske klase $x_i$ pronalazi se $k$ najbli\v{z}ih susjeda unutar iste
klase (Euklidskom udaljeno\v{s}\'{c}u); nasumi\v{c}no se odabire jedan
susjed $x_{nn}$; novi sinteti\v{c}ki primjer generira se linearnom
interpolacijom prema formuli $x_{new} = x_i + \lambda \cdot (x_{nn} - x_i)$,
gdje je $\lambda \sim U(0,1)$. Postupak se ponavlja dok se ne postigne
\v{z}eljeni broj sinteti\v{c}kih primjera. Bit \'{c}e prikazan i
pseudo-k\^{o}d algoritma.

Spomenut \'{c}e se parametri algoritma: broj najbli\v{z}ih susjeda $k$
(uobi\v{c}ajeno $k = 5$ prema preporuci iz izvornog rada) i postotak
preuzorkovanja $N$ (naj\v{c}e\v{s}\'{c}e se cilja na potpuno balansiranje).
Naglasit \'{c}e se da izvorni SMOTE generira isklju\v{c}ivo to\v{c}ke na
linijskim segmentima, \v{s}to ograni\v{c}ava raznolikost sinteti\v{c}kih
uzoraka i motivira razvoj geometrijskih pro\v{s}irenja opisanih u
tre\'{c}em poglavlju.

\subsubsection{Prednosti i nedostaci}

Navest \'{c}e se i kratko obrazlo\v{z}iti prednosti i nedostaci algoritma
SMOTE. Me\dj{}u prednostima istaknut \'{c}e se: smanjenje overfittinga u
odnosu na nasumi\v{c}no preuzorkovanje, pro\v{s}irenje regije odluke
manjinske klase, jednostavnost implementacije, neovisnost o klasifikatoru
te fleksibilnost u pogledu omjera preuzorkovanja.

Me\dj{}u nedostacima obradit \'{c}e se: prekomjerna generalizacija jer
algoritam generira primjere neovisno o distribuciji ve\'{c}inske klase
\cite{stefanowski2007}, osjetljivost na \v{s}um i outlier-e, fiksni broj
susjeda $k$ koji sve primjere tretira jednako, problem s manifold podacima
\cite{bellinger2016}, pote\v{s}ko\'{c}e s opravdavanjem sinteti\v{c}kih
primjera u domenama poput medicine te smanjena u\v{c}inkovitost u
visokodimenzionalnim prostorima. Pregled navedenih ograni\v{c}enja
temeljit \'{c}e se na radu \cite{fernandez2018}.


% ================================================================
% 3. IZVEDENICE ALGORITMA SMOTE I ALTERNATIVNI PRISTUPI
% ================================================================
\section{IZVEDENICE ALGORITMA SMOTE I ALTERNATIVNI PRISTUPI}
% ============================================================
% 03-izvedenice-i-alternative.tex — PRVA VERZIJA
% ============================================================

Temeljni SMOTE algoritam, unato\v{c} svojoj popularnosti, ima nekoliko
uo\v{c}enih nedostataka --- od prekomjerne generalizacije do generiranja
sinteti\v{c}kih primjera u podru\v{c}jima ve\'{c}inske klase
\cite{fernandez2018}. Ovi nedostaci motivirali su razvoj brojnih
izvedenica koje nastoje adresirati pojedine probleme. U ovom poglavlju
bit \'{c}e opisane klju\v{c}ne izvedenice SMOTE algoritma, uklju\v{c}uju\'{c}i
geometrijska pro\v{s}irenja, te alternativni pristupi rje\v{s}avanju
problema neuravnote\v{z}enosti klasa.

\subsection{Borderline-SMOTE}

Borderline-SMOTE \cite{han2005} fokusira se na grani\v{c}ne primjere
manjinske klase koji su klju\v{c}ni za odre\dj{}ivanje granice odluke.
Opisat \'{c}e se kategorizacija manjinskih primjera u tri skupine na
temelju broja ve\'{c}inskih susjeda me\dj{}u $m$ najbli\v{z}ih: sigurni
(Safe), grani\v{c}ni (Danger) i \v{s}um (Noise). Zatim se opisuju
dvije varijante: Borderline-SMOTE1 koja preuzorkuje samo Danger primjere
koriste\'{c}i isklju\v{c}ivo manjinske susjede, i Borderline-SMOTE2 koja
dopu\v{s}ta i ve\'{c}inske susjede uz faktor $\lambda \in (0, 0.5)$.
Dat \'{c}e se pseudo-k\^{o}d za obje varijante. Prednost je fokusiranje
na grani\v{c}ne primjere, a nedostatak osjetljivost na odabir parametra
$m$.

\subsection{ADASYN}

ADASYN (Adaptive Synthetic Sampling) \cite{he2008adasyn} uvodi adaptivni
pristup: vi\v{s}e sinteti\v{c}kih primjera generira se u podru\v{c}jima
gdje je ve\'{c}inska klasa gu\v{s}\'{c}a u susjedstvu manjinskih primjera.
Opisat \'{c}e se ra\v{c}unanje omjera $\Gamma_i = \Delta_i / k$ gdje je
$\Delta_i$ broj ve\'{c}inskih susjeda me\dj{}u $k$ najbli\v{z}ih,
normalizacija u distribuciju $\hat{\Gamma}_i$ te raspodjela ukupnog broja
sinteti\v{c}kih primjera $G$ prema toj distribuciji. Dat \'{c}e se
pseudo-k\^{o}d. Automatsko fokusiranje na te\v{s}ke primjere je glavna
prednost, dok je mana osjetljivost na outlier-e koji dobivaju najvi\v{s}e
sinteti\v{c}kih primjera.

\subsection{Safe-Level-SMOTE}

Safe-Level-SMOTE \cite{bunkhumpornpat2009} uvodi koncept sigurnosne razine
(safe-level) za svaki manjinski primjer --- broj manjinskih susjeda me\dj{}u
$k$ najbli\v{z}ih. Opisat \'{c}e se strategije generiranja ovisno o omjeru
sigurnosnih razina originalnog primjera $sl_p$ i susjeda $sl_n$: ako su
obje niske, generiranje se izbjegava; ako je jedna nula, koristi se kopija
sigurnijeg primjera; ako je $sl_p < sl_n$, novi primjer smje\v{s}ta se
bli\v{z}e sigurnijem. Dat \'{c}e se pseudo-k\^{o}d. Izbjegavanje
generiranja u nesigurnim podru\v{c}jima je prednost, dok je nedostatak
konzervativnost koja mo\v{z}e ograni\v{c}iti pro\v{s}irenje manjinske
klase.

\subsection{K-Means SMOTE}

K-Means SMOTE \cite{douzas2018} rje\v{s}ava problem unutar-klasnog
disbalansa. Postupak: prvo se manjinski primjeri grupiraju K-Means
algoritmom, zatim se identificiraju rijetki klasteri i dodjeljuje im se
proporcionalno vi\v{s}e sinteti\v{c}kih primjera (te\v{z}ina $\propto
1/\text{veli\v{c}ina\_klastera}$), nakon \v{c}ega se SMOTE primjenjuje
unutar svakog klastera koriste\'{c}i samo susjede iz istog klastera.
Dat \'{c}e se pseudo-k\^{o}d. Istovremeno rje\v{s}avanje me\dj{}u-klasnog
i unutar-klasnog disbalansa predstavlja glavnu prednost, dok su dodatni
parametar (broj klastera) i ve\'{c}a ra\v{c}unska zahtjevnost glavni
nedostaci. Srodni pristupi temeljeni na filtriranju obra\dj{}eni su
u radu \cite{saez2015}.

\subsection{SVM-SMOTE}

SVM-SMOTE \cite{nguyen2011} koristi SVM klasifikator za identifikaciju
potpornih vektora manjinske klase koji le\v{z}e na granici odluke.
Postupak: trenira se SVM na originalnim podacima, identificiraju se
potporni vektori manjinske klase, a zatim se generiraju sinteti\v{c}ki
primjeri du\v{z} linija koje spajaju te potporne vektore s njihovim
najbli\v{z}im susjedima. Dat \'{c}e se pseudo-k\^{o}d. Precizno
pozicioniranje novih primjera na granici odluke je klju\v{c}na prednost,
a ve\'{c}a ra\v{c}unska slo\v{z}enost zbog treniranja SVM-a glavni
nedostatak.

\subsection{SMOTE-ENN i SMOTE-Tomek}

SMOTE-ENN i SMOTE-Tomek \cite{batista2004} kombiniraju preuzorkovanje s
naknadnim \v{c}i\v{s}\'{c}enjem skupa podataka. SMOTE-ENN prvo generira
sinteti\v{c}ke primjere, a zatim uklanja sve primjere koje ve\'{c}ina
$k$-NN susjeda pogre\v{s}no klasificira. SMOTE-Tomek tako\dj{}er prvo
primjenjuje SMOTE, a zatim uklanja Tomek Link parove --- parove primjera
razli\v{c}itih klasa koji su me\dj{}usobno najbli\v{z}i. Dat \'{c}e se
pseudo-k\^{o}d za oba pristupa. Prednost je smanjenje \v{s}uma i
preklapanja klasa, dok je nedostatak mogu\'{c}nost gubitka korisnih
primjera.

\subsection{Geometrijska pro\v{s}irenja SMOTE-a}

Izvorni SMOTE generira sinteti\v{c}ke primjere isklju\v{c}ivo na linijskom
segmentu izme\dj{}u dva postoje\'{c}a primjera, \v{s}to ograni\v{c}ava
raznolikost generiranih uzoraka. U nastavku se opisuju tri geometrijska
pro\v{s}irenja koja adresiraju ovo ograni\v{c}enje.

G-SMOTE (Geometric SMOTE) \cite{douzas2019} generira sinteti\v{c}ke
primjere unutar vi\v{s}edimenzionalnog geometrijskog sektora definiranog
manjinskim primjerom kao sredi\v{s}tem i smjerom prema odabranom susjedu.
Uvodi dva dodatna parametra: faktor deformacije i faktor skra\'{c}ivanja
koji kontroliraju koliko daleko od linije se primjer mo\v{z}e generirati.
G-SMOTE se mo\v{z}e koristiti kao izravna zamjena za osnovni SMOTE.

Random SMOTE \cite{dong2011} za svaki manjinski primjer nasumi\v{c}no bira
smjer i udaljenost za generiranje, \v{c}ime pove\'{c}ava raznolikost ali
mo\v{z}e generirati u podru\v{c}jima ve\'{c}inske klase.

Polynom-Fit SMOTE \cite{gazzah2008} koristi polinomnu interpolaciju kroz
vi\v{s}e susjeda umjesto linearnog segmenta, \v{s}to bolje prati stvarnu
geometriju podataka kod nelinearnih distribucija. Jo\v{s} jedno
zna\v{c}ajno geometrijsko/te\v{z}insko pro\v{s}irenje koje nije
implementirano u ovom radu je MWMOTE \cite{barua2014}.

\subsection{Usporedni pregled izvedenica}

Dat \'{c}e se tablica sa sumarnim pregledom svih obra\dj{}enih izvedenica,
uklju\v{c}uju\'{c}i referencu, godinu objave, klju\v{c}ni problem koji
rje\v{s}ava, osnovni koncept i klju\v{c}ne parametre. Dodatne varijante
poput MWMOTE \cite{barua2014} i SMOTE-IPF \cite{saez2015} postoje u
literaturi, no izvan su opsega ove eksperimentalne usporedbe.

\subsection{Alternativni pristupi rje\v{s}avanju problema}

\subsubsection{Metode poduzorkovanja}

Opisat \'{c}e se metode koje smanjuju dominaciju ve\'{c}inske klase
uklanjanjem njenih primjera: NearMiss, Tomek Links, Edited Nearest
Neighbors (ENN) i One-Sided Selection (OSS). Za svaku \'{c}e se objasniti
princip rada i situacije u kojima se koristi. Koristit \'{c}e se
literatura \cite{batista2004,liu2009}.

\subsubsection{Cost-sensitive u\v{c}enje}

Opisat \'{c}e se princip cost-sensitive u\v{c}enja \cite{elkan2001} kod
kojeg se razli\v{c}itim klasama dodjeljuju razli\v{c}ite cijene pogre\v{s}ne
klasifikacije, \v{c}ime se klasifikator prisiljava da posveti vi\v{s}e
pa\v{z}nje manjinskoj klasi. Objasnit \'{c}e se implementacija putem
matrice tro\v{s}kova i class\_weight parametra. Navest \'{c}e se prednosti
(ne mijenja distribuciju podataka) i nedostaci (zahtijeva pa\v{z}ljivo
definiranje tro\v{s}kova).

\subsubsection{Ansambl metode}

Opisat \'{c}e se ansambl pristupi koji kombiniraju resampliranje s
bagging i boosting metodama: SMOTEBoost \cite{chawla2003} primjenjuje
SMOTE u svakoj iteraciji AdaBoost-a, RUSBoost \cite{seiffert2010} koristi
nasumi\v{c}no poduzorkovanje, EasyEnsemble i BalanceCascade \cite{liu2009}
koriste vi\v{s}estruko poduzorkovanje ve\'{c}inske klase, a
BalancedRandomForest balansira bootstrap uzorke. Pregled \'{c}e se
temeljiti na radu \cite{galar2012}.

\subsubsection{Pristupi zasnovani na neuronskim mre\v{z}ama}

Opisat \'{c}e se Focal Loss funkcija \cite{lin2017} koja modificira
cross-entropy smanjuju\'{c}i doprinos lakih primjera, te generativne
suprotstavljene mre\v{z}e \cite{goodfellow2014} za generiranje
visokokvalitetnih sinteti\v{c}kih primjera manjinske klase. Navest
\'{c}e se prednosti (kvaliteta generiranja) i nedostaci (veliki zahtjevi
za koli\v{c}inom podataka i ra\v{c}unskim resursima). Ovaj dio \'{c}e biti
pregledan jer fokus rada ostaje na SMOTE algoritmima.


% ================================================================
% 4. OSTVARENO PROGRAMSKO RJE\v{S}ENJE
% ================================================================
\section{OSTVARENO PROGRAMSKO RJE\v{S}ENJE}
% ============================================================
% 04-programsko-rjesenje.tex - PRVA VERZIJA
% ============================================================

Programsko rje\v{s}enje razvijeno je u programskom jeziku Python, a
obuhva\'{c}a samostalnu implementaciju svih jedanaest SMOTE algoritama,
evaluacijski okvir za eksperimentalnu analizu s osam klasifikatora i
vi\v{s}e metrika evaluacije, modul za statisti\v{c}ku obradu rezultata
te web su\v{c}elje za interaktivnu vizualizaciju. U nastavku se
detaljno opisuju arhitektura i na\v{c}in rada.

\subsection{Na\v{c}in rada programskog rje\v{s}enja}

Arhitektura rje\v{s}enja sastoji se od nekoliko me\dj{}usobno povezanih
modula. Jezgru \v{c}ini paket \texttt{smote\_variants} s apstraktnom
baznom klasom \texttt{BaseSMOTE} koja definira su\v{c}elje
\texttt{fit\_resample(X, y)} \v{s}to ga svaka izvedenica mora
implementirati. Ovo omogu\'{c}uje uniformno kori\v{s}tenje svih
algoritama u evaluacijskom okviru. Implementirane su sljede\'{c}e
varijante: SMOTE, Borderline-SMOTE (obje ina\v{c}ice --- BS1 i BS2),
ADASYN, Safe-Level-SMOTE, K-Means SMOTE, SVM-SMOTE, SMOTE-ENN,
SMOTE-Tomek, Geometric SMOTE (G-SMOTE), Random SMOTE i Polynom-Fit SMOTE.
Svi algoritmi implementirani su samostalno (from-scratch), oslanjaju\'{c}i
se na biblioteke \texttt{numpy} za numeri\v{c}ke operacije i
\texttt{scikit-learn} \cite{pedregosa2011} isklju\v{c}ivo za KNN
pretra\v{z}ivanje (\texttt{NearestNeighbors}), klasteriranje
(\texttt{KMeans}) i SVM (\texttt{SVC}) kod odgovaraju\'{c}ih izvedenica.

Za svaki SMOTE algoritam predvi\dj{}eno je iterativno ispitivanje vi\v{s}e
vrijednosti klju\v{c}nih parametara --- primjerice, broj susjeda $k$
testirat \'{c}e se u rasponu $\{3, 5, 7, 10\}$ --- kako bi se za svaki
skup podataka i klasifikator identificirala konfiguracija koja daje
najbolje rezultate. Ovaj pristup omogu\'{c}uje pravednu usporedbu
algoritama jer svaki od njih dolazi do izra\v{z}aja u svojoj optimalnoj
postavci, \v{c}ime se izbjegava pristranost koja bi nastala fiksiranjem
istih vrijednosti parametara za sve varijante.

Evaluacijski modul implementira stratificiranu $k$-struku unakrsnu
validaciju ($k = 5$) za svaku kombinaciju algoritma preuzorkovanja i
klasifikatora. Kori\v{s}teni klasifikatori uklju\v{c}uju stablo odluke
(DT), slu\v{c}ajnu \v{s}umu (RF), XGBoost \v{c}ija je prednost ugra\dj{}ena
podr\v{s}ka za neuravnote\v{z}ene klase kroz parametar
\texttt{scale\_pos\_weight}, logisti\v{c}ku regresiju (LR), SVM s RBF
jezgrom, $k$-NN, Gaussov naivni Bayes (GNB) te vi\v{s}eslojni perceptron
(MLP) s jednim skrivenim slojem kao jednostavna neuronska referentna
to\v{c}ka. Svi klasifikatori preuzeti su iz biblioteke scikit-learn, uz
iznimku XGBoosta koji dolazi iz istoimene biblioteke.

Za evaluaciju performansi prikupljat \'{c}e se \v{s}iri skup metrika
nego \v{s}to je to uobi\v{c}ajeno u sli\v{c}nim radovima, s ciljem
\v{s}to potpunijeg i znanstveno utemeljenog uvida u pona\v{s}anje
svakog algoritma. Iako su neke od ovih mjera ra\v{c}unski zahtjevnije
od jednostavne to\v{c}nosti klasifikacije --- primjerice, AUC-PR zahtijeva
izra\v{c}un krivulje u vi\v{s}e to\v{c}aka, a MCC uzima u obzir sva
\v{c}etiri elementa matrice zbunjenosti --- upravo ta slo\v{z}enost
omogu\'{c}uje pouzdaniju procjenu na neuravnote\v{z}enim podacima.
Prikupljat \'{c}e se sljede\'{c}e metrike:
F-1 mjera kao harmonijska sredina preciznosti i odziva; G-Mean kao
geometrijska sredina odziva i specifi\v{c}nosti, \v{s}to je posebno
prikladno za neuravnote\v{z}ene podatke jer ka\v{z}njava modele koji
favoriziraju jednu klasu; AUC-ROC kao mjera separabilnosti klasa neovisna
o pragu odluke; AUC-PR (povr\v{s}ina ispod Precision-Recall krivulje)
koja je informativnija od AUC-ROC kod ekstremne neuravnote\v{z}enosti;
Balanced Accuracy kao aritmeti\v{c}ka sredina odziva i specifi\v{c}nosti;
Matthewsov koeficijent korelacije (MCC) koji pru\v{z}a sveobuhvatnu ocjenu
temeljem sva \v{c}etiri elementa matrice zbunjenosti i smatra se jednom od
najpouzdanijih pojedina\v{c}nih mjera za neuravnote\v{z}ene probleme; te
F2-mjera koja daje ve\'{c}u te\v{z}inu odzivu, \v{s}to je relevantno u
domenama gdje su la\v{z}no negativni primjeri skuplji od la\v{z}no
pozitivnih.

Modul za statisti\v{c}ku analizu provodit \'{c}e tri testa \v{c}iji
\'{c}e se rezultati koristiti pri interpretaciji eksperimentalnih
nalaza u petom poglavlju. Friedmanov test \cite{demsar2006} slu\v{z}i za
utvr\dj{}ivanje postoje li statisti\v{c}ki zna\v{c}ajne razlike me\dj{}u
algoritmima na razini svih skupova podataka --- odnosno, nije li uo\v{c}ena
razlika u performansama tek posljedica slu\v{c}ajnosti. U slu\v{c}aju
odbacivanja nul-hipoteze Friedmanovog testa, Nemenyi post-hoc test
identificira koji konkretni parovi algoritama se statisti\v{c}ki
zna\v{c}ajno razlikuju, a rezultati se prikazuju kriti\v{c}nim
dijagramima. Wilcoxon signed-rank test koristit \'{c}e se za direktnu
usporedbu svake izvedenice s osnovnim SMOTE algoritmom kao referentnom
to\v{c}kom, \v{c}ime se dobiva odgovor na pitanje donosi li odre\dj{}ena
modifikacija SMOTE-a zna\v{c}ajno pobolj\v{s}anje ili ne. Svi testovi
provodit \'{c}e se kori\v{s}tenjem biblioteke \texttt{scipy.stats}.

Vizualizacijski modul slu\v{z}i za grafi\v{c}ki prikaz rezultata kroz
CD dijagrame, box-plotove, heatmap-e i scatter prikaze sinteti\v{c}kih
primjera.

\subsection{Prikaz i na\v{c}in uporabe programskog rje\v{s}enja}

Eksperimenti se pokre\'{c}u iz komandne linije naredbom
\texttt{python run\_analysis.py}, pri \v{c}emu je mogu\'{c}e konfigurirati
parametre SMOTE algoritama, odabir klasifikatora i skupova podataka.
Rezultati eksperimenata pohranjuju se u CSV formatu i prikazuju kroz
tablice i grafove. Web su\v{c}elje --- opisano u nastavku ovog poglavlja --- omogu\'{c}it
\'{c}e interaktivnu vizualizaciju procesa generiranja sinteti\v{c}kih
primjera. Za ilustraciju \'{c}e se koristiti screenshotovi glavnih
dijelova su\v{c}elja.

Eventualno, dati dijagram toka programskog rje\v{s}enja na visokoj razini
koji prikazuje tijek od u\v{c}itavanja podataka do prikaza rezultata.
Klju\v{c}ni dijelovi programskog k\^{o}da bit \'{c}e prilo\v{z}eni u
prilogu.

\subsection{Web su\v{c}elje za vizualizaciju}

% ============================================================
% 06-web-sucelje.tex — PRVA VERZIJA
% ============================================================

Razvijeno je i web su\v{c}elje za interaktivnu
vizualizaciju procesa generiranja sinteti\v{c}kih primjera razli\v{c}itim
SMOTE varijantama. Su\v{c}elje omogu\'{c}uje korisniku da na intuitivan
na\v{c}in istra\v{z}i razlike me\dj{}u algoritmima i utjecaj njihovih
parametara na distribuciju generiranih uzoraka.

\subsubsection{Motivacija}

Razumijevanje na\v{c}ina na koji razli\v{c}iti SMOTE algoritmi generiraju
sinteti\v{c}ke primjere \v{c}esto je ote\v{z}ano zbog vi\v{s}edimenzionalne
prirode podataka. Vizualizacija u dvije ili tri dimenzije, dobivena
redukcijom dimenzionalnosti, omogu\'{c}uje intuitivno sagledavanje
razlika izme\dj{}u linearnih pristupa (klasi\v{c}ni SMOTE) i
geometrijskih pro\v{s}irenja (G-SMOTE), \v{s}to ima i zna\v{c}ajnu
edukacijsku vrijednost.

\subsubsection{Arhitektura web su\v{c}elja}

Su\v{c}elje je realizirano kori\v{s}tenjem biblioteke Streamlit za
Python koja omogu\'{c}uje brzo prototipiranje interaktivnih web
aplikacija. Backend se izravno oslanja na implementirane SMOTE algoritme
iz paketa \texttt{smote\_variants}. Za vizualizaciju se koristi Plotly
(interaktivni scatter plotovi), dok se redukcija dimenzionalnosti provodi
algoritmima PCA i t-SNE iz biblioteke scikit-learn.

\subsubsection{Funkcionalnosti}

Opisuju se klju\v{c}ne funkcionalnosti su\v{c}elja:
u\v{c}itavanje vlastitog CSV skupa podataka ili odabir jednog od
ugra\dj{}enih skupova; redukcija dimenzionalnosti na 2D/3D pomo\'{c}u
PCA ili t-SNE; odabir jedne ili vi\v{s}e SMOTE varijanti za usporedbu;
pode\v{s}avanje parametara ($k$, postotak preuzorkovanja); usporedni
prikaz originalnih i sinteti\v{c}kih primjera; animacija procesa
generiranja za 2D podatke; prikaz granice odluke klasifikatora prije i
poslije preuzorkovanja te tablica metrika za kvantitativnu usporedbu.

\subsubsection{Prikazi web su\v{c}elja}

Prikazat \'{c}e se screenshotovi glavnih ekrana su\v{c}elja s kratkim
opisima: po\v{c}etni ekran s kontrolama, PCA prikaz originalnih podataka,
usporedni prikaz vi\v{s}e SMOTE varijanti, prikaz generiranja G-SMOTE u
odnosu na osnovni SMOTE, granica odluke prije i poslije preuzorkovanja
te tablica metrika.



% ================================================================
% 5. EKSPERIMENTALNA ANALIZA
% ================================================================
\section{EKSPERIMENTALNA ANALIZA}
% ============================================================
% 05-eksperimentalna-analiza.tex — PRVA VERZIJA
% ============================================================

Eksperimentalna analiza ima za cilj usporediti u\v{c}inkovitost
jedanaest SMOTE varijanti na skupovima podataka razli\v{c}itih
karakteristika. Analizirat \'{c}e se utjecaj omjera neuravnote\v{z}enosti,
dimenzionalnosti i prisutnosti \v{s}uma na relativne performanse
algoritama, a statisti\v{c}kim testovima provjerit \'{c}e se
zna\v{c}ajnost uo\v{c}enih razlika.

\subsection{Opis skupova podataka}

Opisuju se skupovi podataka kori\v{s}teni u eksperimentalnoj
analizi. Koristit \'{c}e se kombinacija stvarnih i sinteti\v{c}kih
skupova. Za stvarne skupove bit \'{c}e odabran odre\dj{}eni broj
neuravnote\v{z}enih skupova podataka iz razli\v{c}itih domena, pri
\v{c}emu \'{c}e se nastojati obuhvatiti \v{s}to raznovrsniji raspon
omjera neuravnote\v{z}enosti, dimenzionalnosti i veli\v{c}ina kako bi
se dobio \v{s}to op\'{c}enitiji uvid u pona\v{s}anje algoritama.
Tako\dj{}er, poku\v{s}at \'{c}e se identificirati skupovi podataka koji
donekle prate odre\dj{}ene statisti\v{c}ke distribucije (primjerice,
normalnu, eksponencijalnu ili multimodalnu) kako bi se ispitao utjecaj
samih algoritama na podatke razli\v{c}itih distribucijskih svojstava.

Sinteti\v{c}ki skupovi generirat \'{c}e se s kontroliranim parametrima --- omjerom neuravnote\v{z}enosti, dimenzionalno\v{s}\'{c}u i razinom
\v{s}uma --- \v{c}ime \'{c}e se omogu\'{c}iti ciljano ispitivanje
utjecaja pojedinih karakteristika na performanse algoritama, neovisno
o specifi\v{c}nostima stvarnih domena. Osvrnuti se na ograni\v{c}enja
sinteti\v{c}kih skupova podataka.

Dat \'{c}e se tablica s karakteristikama svih kori\v{s}tenih skupova:
naziv, broj primjera, broj zna\v{c}ajki, omjer neuravnote\v{z}enosti
te kategorija prema IR-u (nizak, srednji, visok).

\subsection{Postavke eksperimenta}

Eksperimenti \'{c}e se provesti kori\v{s}tenjem 5-struke stratificirane
unakrsne validacije, pri \v{c}emu je u svakoj iteraciji omjer
neuravnote\v{z}enosti o\v{c}uvan u skupu za treniranje i testiranje. Za
svaki SMOTE algoritam testirat \'{c}e se vrijednosti parametra
$k \in \{3, 5, 7, 10\}$, a preuzorkovanje \'{c}e se vr\v{s}iti do
potpune balansiranosti klasa. Klasifikatori uklju\v{c}uju stablo odluke,
slu\v{c}ajnu \v{s}umu sa 100 stabala, XGBoost, logisti\v{c}ku regresiju,
SVM s RBF jezgrom, $k$-NN s $k=5$, Gaussov naivni Bayes i vi\v{s}eslojni
perceptron (MLP). XGBoost koristi parametar \texttt{scale\_pos\_weight}
za nativno rukovanje neuravnote\v{z}eno\v{s}\'{c}u klasa, dok MLP
s jednim skrivenim slojem od 100 neurona slu\v{z}i kao neuronska
referentna to\v{c}ka. Kao bazne metode koristit
\'{c}e se klasifikacija bez preuzorkovanja, nasumi\v{c}no poduzorkovanje
i nasumi\v{c}no preuzorkovanje.

Prije klasifikacije provest \'{c}e se standardizacija zna\v{c}ajki
pomo\'{c}u \texttt{StandardScaler}-a. Svaka konfiguracija ponovit
\'{c}e se 30 puta s razli\v{c}itim po\v{c}etnim vrijednostima generatora
slu\v{c}ajnih brojeva. Prikupljat \'{c}e se F1, G-Mean, AUC-ROC,
Balanced Accuracy i MCC metrike, a rezultati \'{c}e se prikazati kao
aritmeti\v{c}ka sredina $\pm$ standardna devijacija kroz svih 30
ponavljanja. Navest \'{c}e se i hardverska te softverska konfiguracija
na kojoj su eksperimenti izvedeni.

\subsection{Rezultati}

Zbog velikog broja kombinacija algoritama, klasifikatora i skupova
podataka, rezultati \'{c}e se prikazati na vi\v{s}e komplementarnih
na\v{c}ina, svaki s razli\v{c}itom razinom granularnosti. U glavnom
tekstu fokus \'{c}e biti na agregiranim i statisti\v{c}ki obra\dj{}enim
prikazima koji omogu\'{c}uju uo\v{c}avanje \v{s}irih obrazaca, dok
\'{c}e kompletni rezultati svih eksperimenata biti dani u prilogu.

Prvo, za svaku metriku dat \'{c}e se zbirna tablica s prosje\v{c}nim
rangom svakog algoritma kroz sve skupove podataka i klasifikatore,
poredana od najboljeg prema najlo\v{s}ijem. Ovakav prikaz omogu\'{c}uje
brzi uvid u ukupni poredak algoritama bez zatrpavanja detaljima.
Zatim \'{c}e se za svaku metriku prikazati kriti\v{c}ni dijagram
(CD dijagram) koji vizualno sa\v{z}ima rezultate Friedmanovog i Nemenyi
testa --- algoritmi koji se statisti\v{c}ki zna\v{c}ajno ne razlikuju
povezani su vodoravnom crtom.

Za detaljniji uvid, prikazat \'{c}e se heatmap-ovi koji za svaki par
(algoritam, skup podataka) prikazuju vrijednost odabrane metrike uz
kodiranje bojom, \v{s}to omogu\'{c}uje uo\v{c}avanje obrazaca poput
toga koji algoritmi dominiraju na skupovima s visokim IR-om, a koji na
onima s niskim. Box-plotovi \'{c}e se koristiti za prikaz distribucije
rezultata po algoritmima kroz sva ponavljanja, \v{c}ime se dobiva uvid
ne samo u prosje\v{c}nu performansu ve\'{c} i u njenu varijabilnost.

Za svaki klasifikator izdvojit \'{c}e se i kratko prokomentirati
najva\v{z}niji uo\v{c}eni obrasci --- primjerice, kod kojih klasifikatora
preuzorkovanje donosi najve\'{c}e pobolj\v{s}anje, kako se pona\v{s}aju
geometrijska pro\v{s}irenja u odnosu na klasi\v{c}ne varijante te u kojim
situacijama SMOTE nema u\v{c}inka ili \v{c}ak pogor\v{s}ava rezultate.

Kompletne tablice sa svim vrijednostima metrika za svaku kombinaciju
algoritma, klasifikatora i skupa podataka bit \'{c}e prilo\v{z}ene u
prilogu, organizirane po metrikama radi preglednosti.


\subsection{Statisti\v{c}ka analiza}

Za statisti\v{c}ku usporedbu performansi algoritama koristit \'{c}e se
metodologija Dem\v{s}ara \cite{demsar2006}. Provest \'{c}e se
Friedmanov test za utvr\dj{}ivanje postoje li zna\v{c}ajne razlike
me\dj{}u algoritmima, Nemenyi post-hoc test za identifikaciju razlika
izme\dj{}u konkretnih parova (prikazano kriti\v{c}nim dijagramima) te
Wilcoxon signed-rank test za usporedbu svake izvedenice s osnovnim
SMOTE-om. Dodatno, skupovi \'{c}e se grupirati prema omjeru
neuravnote\v{z}enosti, dimenzionalnosti i razini \v{s}uma kako bi se
ispitalo kako ovi faktori utje\v{c}u na relativne performanse
algoritama.

\subsection{Diskusija}

Na temelju provedene eksperimentalne i statisti\v{c}ke analize
interpretirat \'{c}e se dobiveni rezultati s naglaskom na sljede\'{c}a
pitanja: koji SMOTE algoritam pokazuje najbolje ukupne performanse,
u kojim situacijama se preporu\v{c}uje koja izvedenica, kako
geometrijska pro\v{s}irenja stoje u odnosu na klasi\v{c}ne varijante,
kakav je utjecaj parametra $k$ na razli\v{c}ite klasifikatore te koja
je ra\v{c}unska zahtjevnost pojedinih algoritama. Dani \'{c}e se i
prakti\v{c}ne preporuke za odabir algoritma temeljem karakteristika
skupa podataka. Rezultati \'{c}e se usporediti s nalazima iz literature
\cite{bajer2019,dudjak2020,kovacs2019,haixiang2017}.


% ================================================================
% 6. ZAKLJU\v{C}AK
% ================================================================
\section{ZAKLJU\v{C}AK}
% ============================================================
% 07-zakljucak.tex — PRVA VERZIJA
% ============================================================

U zaklju\v{c}ku dajemo osvrt na cjelokupni rad, sa\v{z}eto navodimo
\v{s}to je napravljeno i \v{s}to su klju\v{c}ni nalazi, a zatim
navodimo ograni\v{c}enja i otvaramo smjernice za budu\'{c}i rad.

Kroz rad je najprije obra\dj{}en teorijski okvir: problem
neuravnote\v{z}enosti klasa, temeljni SMOTE algoritam i pregled njegovih
jedanaest izvedenica. Zatim je opisano ostvareno programsko rje\v{s}enje
— Python implementacija svih varijanti, evaluacijski okvir s osam
klasifikatora i vi\v{s}e metrika, te modul za statisti\v{c}ku obradu
prema Dem\v{s}arovoj metodologiji. Provedena je eksperimentalna analiza
na stvarnim i sinteti\v{c}kim skupovima razli\v{c}itih karakteristika, a
rezultati su statisti\v{c}ki obra\dj{}eni i interpretirani. Kao dodatak,
izra\dj{}eno je web su\v{c}elje za vizualizaciju.

Klju\v{c}ni nalazi koje o\v{c}ekujemo: identificirat \'{c}emo koji se
SMOTE algoritmi najbolje pokazuju u pojedinim uvjetima, koliko
geometrijska pro\v{s}irenja zaista doprinose u odnosu na osnovni SMOTE,
te kakvu ulogu igra izbor klasifikatora.

Ograni\v{c}enja rada: fokus je na binarnoj klasifikaciji, \v{s}to
zna\v{c}i da zaklju\v{c}ci nisu izravno primjenjivi na vi\v{s}eklasne
probleme. Tako\dj{}er, broj testiranih vrijednosti parametara je
ograni\v{c}en, a sinteti\v{c}ki skupovi ne mogu u potpunosti zamijeniti
stvarne podatke.

Za budu\'{c}i rad otvaramo nekoliko pravaca: pro\v{s}irenje na
vi\v{s}eklasne probleme, dublja integracija s deep learning pristupima,
automatska optimizacija hiperparametara SMOTE algoritama te evaluacija
na ekstremno neuravnote\v{z}enim skupovima (IR $>$ 100) koji su u praksi
sve \v{c}e\v{s}\'{c}i.


% ================================================================
% LITERATURA
% ================================================================

\newpage
\bibliographystyle{unsrt}
\bibliography{references}

% ================================================================
% SA\v{Z}ETAK
% ================================================================

\newpage
\section*{Sa\v{z}etak}
\addcontentsline{toc}{section}{Sa\v{z}etak}
% ============================================================
% 00-sazetak-hr.tex
% ============================================================

Problem neuravnote\v{z}enosti klasa jedan je od klju\v{c}nih izazova u
strojinom u\v{c}enju, posebice u klasifikacijskim zadacima gdje
distribucija primjera me\dj{}u klasama nije uniformna. Algoritam SMOTE
(engl.\ Synthetic Minority Over-sampling Technique) predstavlja jedan od
najpopularnijih pristupa za rje\v{s}avanje ovog problema — generira
sinteti\v{c}ke primjere manjinske klase linearnom interpolacijom
postoje\'{c}ih.

U ovom radu analizirani su temeljni SMOTE algoritam i njegove izvedenice:
Borderline-SMOTE, ADASYN, Safe-Level-SMOTE, K-Means SMOTE, SVM-SMOTE,
kombinirani pristupi SMOTE-ENN i SMOTE-Tomek te geometrijska
pro\v{s}irenja G-SMOTE, Random SMOTE i Polynom-Fit SMOTE. Usporedno su
razmatrani i alternativni pristupi poput cost-sensitive u\v{c}enja,
ansambl metoda i primjene neuronskih mre\v{z}a.

Razvijeno programsko rje\v{s}enje implementira navedene algoritme u
Pythonu i omogu\'{c}uje njihovu eksperimentalnu evaluaciju na \v{s}irokom
spektru skupova podataka razli\v{c}itih karakteristika. Provedena je
statisti\v{c}ka analiza rezultata s ciljem utvr\dj{}ivanja pona\v{s}anja
pojedinih algoritama pri razli\v{c}itim omjerima klasa,
dimenzionalnostima i razinama \v{s}uma. Dodatno, izra\dj{}eno je web
su\v{c}elje za vizualizaciju procesa generiranja sinteti\v{c}kih
primjera.

\bigskip
\noindent\textbf{Klju\v{c}ne rije\v{c}i:} klasifikacija, neuravnote\v{z}enost klasa,
preuzorkovanje, sinteti\v{c}ki uzorci, SMOTE


\newpage
\section*{Abstract}
\addcontentsline{toc}{section}{Abstract}
% ============================================================
% 00-sazetak-en.tex
% ============================================================

The class imbalance problem is one of the key challenges in machine
learning, particularly in classification tasks where the distribution of
examples across classes is not uniform. The SMOTE (Synthetic Minority
Over-sampling Technique) algorithm represents one of the most popular
approaches to address this problem by generating synthetic examples of
the minority class through linear interpolation.

This thesis analyzes the basic SMOTE algorithm and its variants:
Borderline-SMOTE, ADASYN, Safe-Level-SMOTE, K-Means SMOTE, SVM-SMOTE,
the combined approaches SMOTE-ENN and SMOTE-Tomek, and geometric
extensions G-SMOTE, Random SMOTE and Polynom-Fit SMOTE. Alternative
approaches such as cost-sensitive learning, ensemble methods, and neural
network applications are also considered.

The developed software solution implements the described algorithms in
Python and enables their experimental evaluation on a wide range of
datasets with varying characteristics. Statistical analysis of the
results was conducted to determine the behavior of individual algorithms
under different class ratios, dimensionalities, and noise levels.
Additionally, a web interface was developed for visualizing the synthetic
sample generation process.

\bigskip
\noindent\textbf{Keywords:} class imbalance, classification, oversampling,
SMOTE, synthetic samples


% ================================================================
% PRILOZI
% ================================================================

\newpage
\section*{Prilog 1 --- Pro\v{s}ireni rezultati}
\addcontentsline{toc}{section}{Prilog 1 --- Pro\v{s}ireni rezultati}
% ============================================================
% 09-prilozi.tex — PRVA VERZIJA
% ============================================================

U prilozima dajemo materijale koji nisu neophodni za pra\'{c}enje
glavnog teksta, ali pru\v{z}aju dodatne informacije za zainteresiranog
\v{c}itatelja.

\textbf{Prilog P.1.} Pro\v{s}ireni eksperimentalni rezultati: kompletne
tablice F1, G-Mean i MCC vrijednosti za svih osam klasifikatora, za sve
skupove podataka i sve SMOTE varijante. Ovdje su i tablice
$p$-vrijednosti Friedmanovog, Nemenyi i Wilcoxon testova.

\textbf{Prilog P.2.} Klju\v{c}ni dijelovi programskog k\^{o}da:
apstraktna bazna klasa \texttt{BaseSMOTE}, implementacija osnovnog SMOTE
algoritma, te po jedna reprezentativna klasi\v{c}na (Borderline-SMOTE) i
geometrijska (G-SMOTE) izvedenica. Cijeli k\^{o}d je dostupan na
GitHub repozitoriju (URL \'{c}e biti naveden naknadno).

\textbf{Prilog P.3.} Upute za instalaciju ovisnosti i pokretanje:
\begin{verbatim}
pip install -r requirements.txt
python run_analysis.py
streamlit run app.py
\end{verbatim}


\end{document}
